\documentclass[11pt,a4paper]{article}

% ── Typography & Encoding ─────────────────────────────────────────────
\usepackage[utf8]{inputenc}
\usepackage[T1]{fontenc}
\usepackage{newtxtext,newtxmath} % Professional academic font (Times-like)
\usepackage{microtype}           % Advanced typography (kerning, spacing)
\usepackage[margin=1in]{geometry}
\usepackage{setspace}            % Line spacing

% ── Math & Algorithms ─────────────────────────────────────────────────
\usepackage{amsmath}
% amssymb and amsthm are often superseded or clash with newtxmath,
% but amsthm is safe and useful.
\let\openbox\relax
\usepackage{amsthm}
\usepackage{algorithm}
\usepackage[noend]{algpseudocode}

% ── Graphics & Tables ─────────────────────────────────────────────────
\usepackage{graphicx}
\usepackage{booktabs}
\usepackage{multirow}
\usepackage{tabularx}
\usepackage{array}
\usepackage{caption}
\usepackage{subcaption}

% ── Utilities ─────────────────────────────────────────────────────────
\usepackage{xcolor}
\usepackage{enumitem}
\usepackage{siunitx}             % Standardized units and number formatting
\usepackage{authblk}             % Professional author blocks

% ── Hyperref & Cleveref ───────────────────────────────────────────────
\usepackage[colorlinks=true,linkcolor=darkblue,citecolor=darkblue,urlcolor=darkblue]{hyperref}
\usepackage[capitalize,nameinlink]{cleveref}

% ── Formatting adjustments ────────────────────────────────────────────
\definecolor{darkblue}{rgb}{0.0, 0.0, 0.55}
\definecolor{baseline}{HTML}{7f7f7f}
\definecolor{diverse}{HTML}{e377c2}
\definecolor{scalarized}{HTML}{8c564b}

\captionsetup{labelfont=bf, textfont=it, margin=1em}
\setlist{nosep, leftmargin=*}

% ── Algorithmic customisation ─────────────────────────────────────────
\algnewcommand\algorithmicforeach{\textbf{for each}}
\algdef{S}[FOR]{ForEach}[1]{\algorithmicforeach\ #1\ \algorithmicdo}
\algnewcommand\algorithmicinput{\textbf{Input:}}
\algnewcommand\algorithmicoutput{\textbf{Output:}}
\algnewcommand\Input{\item[\algorithmicinput]}
\algnewcommand\Output{\item[\algorithmicoutput]}

% ── Math macros ───────────────────────────────────────────────────────

\newcommand{\calA}{\ensuremath{\mathcal{A}}}
\newcommand{\calE}{\ensuremath{\mathcal{E}}}
\newcommand{\calF}{\ensuremath{\mathcal{F}}}
\newcommand{\calI}{\ensuremath{\mathcal{I}}}
\newcommand{\calN}{\ensuremath{\mathcal{N}}}
\newcommand{\calP}{\ensuremath{\mathcal{P}}}
\newcommand{\calS}{\ensuremath{\mathcal{S}}}
\newcommand{\bbR}{\ensuremath{\mathbb{R}}}
\DeclareMathOperator*{\argmin}{arg\,min}
\DeclareMathOperator*{\argmax}{arg\,max}

% ══════════════════════════════════════════════════════════════════════
\title{%
  \vspace{-1.5cm}
  \textbf{Enhanced Pareto Local Search for the\\
  Satellite Image Mosaic Selection Problem:\\
  Diversity-Guided and Scalarization-Guided\\
  Parent Selection Strategies}
}

\author[1]{Anonymous Author 1}
\author[2]{Anonymous Author 2}
\affil[1]{Department of Computer Science, Anonymous University}
\affil[2]{Institute for Advanced Optimization, Anonymous Institute}
\date{\today}

\begin{document}
\maketitle

% ══════════════════════════════════════════════════════════════════════
\begin{abstract}
\noindent
Pareto Local Search (PLS) is a natural and effective local-search
framework for multi-objective combinatorial optimisation, but its
performance depends critically on how search effort is distributed over
the current approximation set.  In the four-objective Satellite Image
Mosaic Selection (SIMS) problem, a naive PLS policy may spend excessive
time exploring densely populated regions of the archive while neglecting
other regions that are either structurally underrepresented or more
promising under specific trade-off directions.

We study two archive-steering enhancements to a common improved PLS
baseline.  The first is a \emph{diversity-guided probing} strategy that
selects a well-spread subset of candidate parents in objective space
before neighbourhood exploration.  The second is a
\emph{scalarization-guided selection} strategy that repeatedly samples
weighted Chebycheff trade-off directions and explores archive solutions
that are best aligned with those directions.  Both mechanisms operate at
the level of parent-solution scheduling rather than neighbourhood
definition, and therefore complement the standard PLS framework without
changing the underlying local move operator.

Experiments on 20 benchmark instances (five cities, four size tiers,
30--150 images, four objectives) show that both enhancements improve
hypervolume over the baseline on medium and large instances, but that
diversity-guided probing is the stronger default strategy overall.
Across all 20 instances, diversity-guided probing improves final
normalised hypervolume by \(+3.60\%\) on average over the baseline,
whereas scalarization-guided selection improves it by \(+3.42\%\).
Direct comparison shows that diversity-guided probing outperforms
scalarization-guided selection on 18 of 20 instances, with an average
advantage of \(0.18\%\).  Scalarization-guided selection nevertheless
achieves the best result on the hardest Lagos instance
(\(+3.06\%\) over diversity-guided probing at 145 images), indicating
that directional archive steering can be particularly effective on
certain large, structurally heterogeneous instances.

\vspace{0.5em}
\noindent\textbf{Keywords:}\quad multi-objective optimisation, Pareto local search, satellite image selection, hypervolume, diversity-guided search, scalarization-guided search, weighted Chebycheff
\end{abstract}

% ══════════════════════════════════════════════════════════════════════
\section{Introduction}\label{sec:intro}

The \emph{Satellite Image Mosaic Selection} (SIMS) problem requires
choosing a subset of satellite images that jointly cover a region of
interest while balancing several conflicting criteria.  In the
four-objective setting considered here, the optimisation simultaneously
minimises acquisition cost, cloud contamination, worst-case incidence
angle, and resolution loss.  The resulting problem is a multi-objective
variant of weighted set cover and is computationally challenging even
for moderate catalogue sizes.

Pareto Local Search (PLS) is a natural heuristic for this setting.  It
maintains an archive of mutually non-dominated solutions and repeatedly
explores neighbourhoods of selected archive members.  The quality of the
final approximation set depends not only on the neighbourhood operator
itself, but also on the \emph{policy used to decide which archive
solutions should be explored next}.  In many-objective settings, this
policy becomes especially important because the archive can grow rapidly,
making exhaustive and uniformly distributed exploration increasingly
inefficient.

This paper focuses on two search-allocation enhancements that operate at
the parent-selection level:
\begin{enumerate}
  \item \textbf{Diversity-guided probing}, which selects a well-spread
    subset of archive or population members in objective space before
    neighbourhood exploration.
  \item \textbf{Scalarization-guided selection}, which samples weighted
    trade-off directions and prioritises archive solutions that minimise
    the corresponding weighted Chebycheff scalarizations.
\end{enumerate}

Both strategies are layered on top of the same improved PLS baseline and
therefore isolate the effect of \emph{search steering} rather than
neighbourhood design.  The goal is not to alter the definition of local
moves, but to improve the \emph{allocation of limited search effort}
over the current approximation set.

\paragraph{Contributions.}
This paper makes the following contributions:
\begin{enumerate}
  \item We formalise two parent-selection strategies for PLS at the same
    level of abstraction as the baseline algorithm: a diversity-guided
    strategy and a scalarization-guided strategy.
  \item We provide a controlled empirical comparison of these strategies
    against the same baseline PLS on 20 four-objective SIMS instances.
  \item We show that diversity-guided probing is the stronger default
    strategy overall, while scalarization-guided selection can dominate
    on selected large instances.
  \item We analyse the results by instance size and identify the regimes
    in which each strategy is most effective.
\end{enumerate}

% ══════════════════════════════════════════════════════════════════════
\section{Problem Formulation}\label{sec:problem}

We consider a finite set of candidate images
\(\calI = \{1,\dots,m\}\) and a finite set of ground elements
\(\calE = \{1,\dots,n\}\) that must be covered.  Each image
\(i \in \calI\) covers a subset \(C_i \subseteq \calE\).  A feasible
solution is any subset \(S \subseteq \calI\) such that
\[
\bigcup_{i \in S} C_i = \calE.
\]

\subsection{Decision Space}

A solution is represented by a binary vector
\(x \in \{0,1\}^m\), where \(x_i = 1\) if image \(i\) is selected and
\(x_i = 0\) otherwise.  The feasible set is
\[
\calF = \left\{x \in \{0,1\}^m \;\middle|\; \forall e \in \calE,\;
\exists i \in \calI \text{ such that } x_i = 1 \text{ and } e \in C_i
\right\}.
\]

\subsection{Objective Space}

Each feasible solution is evaluated by four minimisation objectives:
\[
f(x) = \bigl(f_1(x), f_2(x), f_3(x), f_4(x)\bigr),
\]
where:
\begin{itemize}
  \item \(f_1\): total acquisition cost,
  \item \(f_2\): total cloud contamination,
  \item \(f_3\): worst-case incidence angle,
  \item \(f_4\): resolution loss.
\end{itemize}

For two feasible solutions \(x,y \in \calF\), we say that \(x\)
\emph{dominates} \(y\), written \(x \prec y\), if
\[
\forall j \in \{1,\dots,4\},\; f_j(x) \le f_j(y)
\quad\text{and}\quad
\exists j \in \{1,\dots,4\},\; f_j(x) < f_j(y).
\]
The goal is to approximate the Pareto set of mutually non-dominated
feasible solutions.

% ══════════════════════════════════════════════════════════════════════
\section{Baseline PLS}\label{sec:baseline}

PLS maintains two sets:
\begin{itemize}
  \item an archive \(\calA\) of mutually non-dominated solutions found
    so far, and
  \item a working population \(\calP \subseteq \calA\) of solutions
    whose neighbourhoods are scheduled for exploration.
\end{itemize}

At each iteration, the algorithm selects one or more parent solutions
from \(\calP\), explores their neighbourhoods, inserts newly discovered
non-dominated neighbours into \(\calA\), and forms the next working
population from the accepted offspring.  The baseline policy considered
here uses the same neighbourhood operator and archive update mechanism
throughout all experiments; only the \emph{parent-selection strategy}
changes.

\subsection{Neighbourhood Operator}

Let \(x \in \calF\) be a feasible solution.  Its neighbourhood
\(\calN(x)\) consists of feasible solutions obtained by local image-set
modifications that preserve full coverage.  The exact move operator is
not the focus of this paper; what matters is that all compared
configurations use the same neighbourhood definition and therefore
differ only in how parent solutions are chosen for exploration.

\subsection{Main Loop}

At a high level, baseline PLS can be written as follows.

\begin{algorithm}[htbp]
\caption{Baseline Pareto Local Search}
\label{alg:pls}
\begin{algorithmic}[1]
\Input Initial feasible population \(\calP_0\), neighbourhood operator \(\calN\)
\Output Approximation archive \(\calA\)
\State \(\calA \gets \textsc{NonDominated}(\calP_0)\)
\State \(\calP \gets \calA\)
\While{termination criterion not met}
  \State choose parent subset \(\calS \subseteq \calP\)
  \State \(\calP' \gets \emptyset\)
  \ForEach{\(x \in \calS\)}
    \ForEach{\(y \in \calN(x)\)}
      \If{\(y\) is not dominated by \(\calA\)}
        \State insert \(y\) into \(\calA\), removing dominated archive members
        \State insert \(y\) into \(\calP'\)
      \EndIf
    \EndFor
  \EndFor
  \State \(\calP \gets \textsc{NextPopulation}(\calP,\calP',\calA)\)
\EndWhile
\State \Return \calA
\end{algorithmic}
\end{algorithm}

The central question addressed in this paper is therefore: \emph{how should the parent subset \(\calS\) be chosen?}

% ══════════════════════════════════════════════════════════════════════
\section{Enhancement 1: Diversity-Guided Probing}\label{sec:diverse}

\subsection{Motivation}

A straightforward parent-selection policy may repeatedly explore archive
members that are close to one another in objective space.  This can
cause search effort to concentrate on already dense regions of the
current approximation set while leaving other regions underexplored.
When the archive is large, such concentration reduces the effective
coverage of the Pareto front approximation.

The purpose of diversity-guided probing is to allocate exploration
effort to a subset of parent solutions that is \emph{well spread} in
objective space.  The underlying principle is that, under a fixed
computational budget, it is often preferable to explore a representative
set of structurally different archive members rather than many nearly
redundant ones.

\subsection{Method}

Let \(\calP\) denote the current candidate parent set.  Diversity-guided
probing selects a subset \(\calS \subseteq \calP\) of prescribed budget
\(B\) such that the selected solutions are as far apart as possible in
normalised objective space.

The method proceeds conceptually in three stages:
\begin{enumerate}
  \item \textbf{Normalisation.}  Objective vectors are mapped to a
    common scale using the current objective ranges in the candidate
    set.
  \item \textbf{Seed selection.}  One candidate is chosen as an initial
    representative.
  \item \textbf{Farthest-point expansion.}  Repeatedly add the candidate
    whose minimum distance to the already selected set is maximal.
\end{enumerate}

This produces a subset that approximates a max--min spread criterion.
The selected parents are then explored in the order in which they were
chosen, so that even if the run is interrupted early, the explored
subset remains broadly representative.

\subsection{Theoretical Effect}

Diversity-guided probing does not change the archive update rule or the
definition of local optimality.  Instead, it changes the \emph{order and
distribution} of neighbourhood exploration.  Its expected effect is:
\begin{itemize}
  \item broader coverage of the current approximation set,
  \item reduced oversampling of dense regions,
  \item improved anytime behaviour under fixed budgets.
\end{itemize}

The method is especially attractive when the archive is already large
and the main challenge is not finding \emph{any} improving move, but
deciding \emph{where} to spend the next unit of search effort.

% ══════════════════════════════════════════════════════════════════════
\section{Enhancement 2: Scalarization-Guided Parent Selection}\label{sec:scalarized}

\subsection{Motivation}

Diversity alone does not distinguish between archive members that are
merely different and archive members that are particularly promising
under a given trade-off direction.  In many-objective search, a useful
alternative is to steer exploration by repeatedly sampling preference
directions and selecting archive solutions that are best aligned with
those directions.

The scalarization-guided strategy follows this idea.  Rather than
choosing parents solely for geometric spread, it samples weighted
trade-off vectors and prioritises archive members that minimise the
corresponding scalarized objective.  This biases the search toward
solutions that are likely to improve the approximation set in specific
regions of the Pareto front.

\subsection{Method}

Let \(z^\star\) denote a reference ideal point and let
\(\lambda = (\lambda_1,\dots,\lambda_D)\) be a non-negative weight
vector with \(\sum_j \lambda_j = 1\).  For a solution \(x\), define the
weighted Chebycheff scalarization
\[
g_\lambda(x \mid z^\star)
=
\max_{j=1,\dots,D}
\left\{
\lambda_j \,\frac{f_j(x)-z^\star_j}{r_j}
\right\},
\]
where \(r_j > 0\) is a normalising range for objective \(j\).

At each parent-selection step:
\begin{enumerate}
  \item sample one or more weight vectors \(\lambda\),
  \item for each sampled \(\lambda\), identify the archive solution
    \(x^\lambda \in \calA\) minimising \(g_\lambda(x \mid z^\star)\),
  \item explore the neighbourhoods of the selected minimisers.
\end{enumerate}

This strategy is archive-driven: the parent candidates are drawn from
the current non-dominated archive rather than from an arbitrary working
population.  In this way, the search is repeatedly redirected toward
solutions that are most representative of sampled trade-off directions.

\subsection{Theoretical Effect}

Scalarization-guided selection changes the search policy from
\emph{coverage-oriented} to \emph{direction-oriented}.  Its expected
effect is:
\begin{itemize}
  \item stronger exploitation of promising trade-off directions,
  \item improved progress in regions where the archive is sparse but
    strategically important,
  \item better alignment between parent selection and quality indicators
    based on directional trade-offs.
\end{itemize}

The method is particularly appealing when the archive is large and the
main challenge is to identify which archive members are most likely to
yield useful improvements under limited exploration budgets.

% ══════════════════════════════════════════════════════════════════════
\section{Experimental Setup}\label{sec:experiments}

\subsection{Instances}

We use 20 benchmark instances derived from real satellite catalogues
over five cities (Lagos, Mexico City, Paris, Rio de Janeiro, and Tokyo
Bay) at four size tiers up to 150 images.  Lagos uses 145 images at the
largest tier, while the other cities use 150 images.  All instances are
optimised with \(D=4\) objectives.

\begin{table}[htbp]
\centering
\caption{Benchmark instance characteristics.}
\label{tab:instances}
\begin{tabular}{lrrr}
\toprule
\textbf{City} & \textbf{Images} & \textbf{Elements} & \textbf{Timeout (s)} \\
\midrule
\multirow{4}{*}{Lagos}
  & 30  &    443  & 300 \\
  & 50  &  1\,147 & 300 \\
  & 100 &  4\,420 & 1200 \\
  & 145 &  8\,503 & 1200 \\
\midrule
\multirow{4}{*}{Mexico City}
  & 30  &    333  & 300 \\
  & 50  &    842  & 300 \\
  & 100 &  3\,384 & 1200 \\
  & 150 &  5\,236 & 1200 \\
\midrule
\multirow{4}{*}{Paris}
  & 30  &    475  & 300 \\
  & 50  &    970  & 300 \\
  & 100 &  4\,503 & 1200 \\
  & 150 & 11\,027 & 1200 \\
\midrule
\multirow{4}{*}{Rio de Janeiro}
  & 30  &    351  & 300 \\
  & 50  &  1\,238 & 300 \\
  & 100 &  5\,905 & 1200 \\
  & 150 & 12\,879 & 1200 \\
\midrule
\multirow{4}{*}{Tokyo Bay}
  & 30  &    298  & 300 \\
  & 50  &    806  & 300 \\
  & 100 &  3\,278 & 1200 \\
  & 150 &  8\,079 & 1200 \\
\bottomrule
\end{tabular}
\end{table}

\subsection{Configurations}

We compare three configurations:
\begin{enumerate}
  \item \textbf{Pure PLS:} the baseline PLS reference.
  \item \textbf{Diverse Probe PLS:} the same baseline augmented with
    diversity-guided parent probing.
  \item \textbf{Scalarized PLS:} the same baseline augmented with
    scalarization-guided archive selection.
\end{enumerate}

The purpose of this design is to isolate the effect of the two
parent-selection strategies while keeping the underlying local search
framework fixed.

\subsection{Metrics}

The primary metric is the \textbf{normalised hypervolume} (HV) computed
with shared per-instance bounds derived from the union of all trace
points across the compared configurations.  This ensures that HV values
are directly comparable.

We also analyse \textbf{HV-over-time} curves reconstructed from the
search traces.  For each configuration and instance, the trace records
the sequence of non-dominated-at-discovery solutions together with their
timestamps.  We reconstruct the Pareto front at evenly spaced time
points and compute HV at each point.

\subsection{Implementation}

The experiment driver is implemented in Python and calls a compiled PLS
solver through language bindings.  Hypervolume-over-time curves are
computed from archived traces using a shared reconstruction and
evaluation pipeline.  All compared configurations use the same
implementation framework and differ only in the parent-selection policy.

% ══════════════════════════════════════════════════════════════════════
\section{Results}\label{sec:results}

\subsection{Hypervolume Summary}

\Cref{tab:results} reports the final normalised hypervolume for all 20
instances.  The table includes the baseline, the diversity-guided
variant, and the scalarization-guided variant.

\begin{table}[htbp]
\centering
\caption{Final normalised hypervolume for Pure PLS, Diverse Probe PLS,
and Scalarized PLS.  Positive \(\Delta\) values indicate improvement
over the Pure PLS baseline.}
\label{tab:results}
\begin{tabular}{lrrrrr}
\toprule
\textbf{Instance} & \textbf{Size} & \textbf{Pure PLS} &
\textbf{Diverse Probe} & \textbf{Scalarized} &
\textbf{Best} \\
\midrule
Lagos 30        &  30 & 0.811595 & 0.811595 & 0.811576 & Diverse \\
Lagos 50        &  50 & 0.868126 & 0.867881 & 0.867573 & Pure \\
Lagos 100       & 100 & 0.886549 & 0.923981 & 0.922884 & Diverse \\
Lagos 145       & 145 & 0.734004 & 0.792120 & 0.816380 & Scalarized \\
\midrule
Mexico City 30  &  30 & 0.942413 & 0.942413 & 0.942400 & Diverse \\
Mexico City 50  &  50 & 0.923096 & 0.922876 & 0.921505 & Pure \\
Mexico City 100 & 100 & 0.902607 & 0.924284 & 0.921899 & Diverse \\
Mexico City 150 & 150 & 0.784194 & 0.883515 & 0.879431 & Diverse \\
\midrule
Paris 30        &  30 & 0.839079 & 0.838858 & 0.821789 & Pure \\
Paris 50        &  50 & 0.853323 & 0.855010 & 0.852591 & Diverse \\
Paris 100       & 100 & 0.782960 & 0.894449 & 0.880452 & Diverse \\
Paris 150       & 150 & 0.875556 & 0.944785 & 0.943737 & Diverse \\
\midrule
Rio 30          &  30 & 0.812803 & 0.812803 & 0.812541 & Diverse \\
Rio 50          &  50 & 0.643427 & 0.643056 & 0.642872 & Pure \\
Rio 100         & 100 & 0.936971 & 0.938717 & 0.937587 & Diverse \\
Rio 150         & 150 & 0.781738 & 0.876039 & 0.876232 & Scalarized \\
\midrule
Tokyo Bay 30    &  30 & 0.719423 & 0.719423 & 0.719303 & Diverse \\
Tokyo Bay 50    &  50 & 0.848188 & 0.848105 & 0.839300 & Pure \\
Tokyo Bay 100   & 100 & 0.880687 & 0.898726 & 0.897927 & Diverse \\
Tokyo Bay 150   & 150 & 0.860256 & 0.932349 & 0.928658 & Diverse \\
\bottomrule
\end{tabular}
\end{table}

Several patterns emerge immediately:
\begin{itemize}
  \item Both enhanced strategies substantially outperform the baseline on
    many medium and large instances.
  \item Diversity-guided probing is the strongest overall strategy,
    achieving the best HV on 18 of 20 instances.
  \item Scalarization-guided selection nevertheless achieves the best
    result on two large instances: Lagos 145 and Rio 150.
  \item On the smallest instances, all three configurations are often
    very close, indicating that the baseline already explores the search
    space effectively under the available budget.
\end{itemize}

\subsection{Improvement over the Baseline}

To quantify the gains over the baseline, \Cref{tab:improvement} reports
the percentage improvement of each enhanced strategy relative to Pure
PLS.

\begin{table}[htbp]
\centering
\caption{Percentage hypervolume improvement over Pure PLS.}
\label{tab:improvement}
\begin{tabular}{lrrr}
\toprule
\textbf{Instance} & \textbf{Size} &
\textbf{Diverse vs Pure (\%)} &
\textbf{Scalarized vs Pure (\%)} \\
\midrule
Lagos 30        &  30 &  +0.00 &  -0.00 \\
Lagos 50        &  50 &  -0.03 &  -0.06 \\
Lagos 100       & 100 &  +4.22 &  +4.10 \\
Lagos 145       & 145 &  +7.92 & +11.22 \\
\midrule
Mexico City 30  &  30 &  +0.00 &  -0.00 \\
Mexico City 50  &  50 &  -0.02 &  -0.17 \\
Mexico City 100 & 100 &  +2.40 &  +2.14 \\
Mexico City 150 & 150 & +12.67 & +12.14 \\
\midrule
Paris 30        &  30 &  -0.03 &  -2.06 \\
Paris 50        &  50 &  +0.20 &  -0.09 \\
Paris 100       & 100 & +14.24 & +12.45 \\
Paris 150       & 150 &  +7.91 &  +7.79 \\
\midrule
Rio 30          &  30 &  +0.00 &  -0.03 \\
Rio 50          &  50 &  -0.06 &  -0.09 \\
Rio 100         & 100 &  +0.19 &  +0.07 \\
Rio 150         & 150 & +12.06 & +12.09 \\
\midrule
Tokyo Bay 30    &  30 &  +0.00 &  -0.02 \\
Tokyo Bay 50    &  50 &  -0.01 &  -1.05 \\
Tokyo Bay 100   & 100 &  +2.05 &  +1.96 \\
Tokyo Bay 150   & 150 &  +8.38 &  +7.95 \\
\midrule
\multicolumn{2}{l}{\textbf{Average}} & \textbf{+3.60} & \textbf{+3.42} \\
\bottomrule
\end{tabular}
\end{table}

The average improvements show that both strategies are beneficial, but
that diversity-guided probing has a slight overall advantage.  The
difference is small in aggregate, yet consistent enough to make
diversity-guided probing the stronger default choice.

\subsection{Direct Comparison: Diverse vs Scalarized}

\Cref{tab:direct} compares the two enhanced strategies directly.

\begin{table}[htbp]
\centering
\caption{Direct comparison of Scalarized PLS against Diverse Probe PLS.
Positive \(\Delta\) indicates that Scalarized PLS is better.}
\label{tab:direct}
\begin{tabular}{lrrr}
\toprule
\textbf{Instance} & \textbf{Size} &
\textbf{Scalarized -- Diverse (\%)} &
\textbf{Winner} \\
\midrule
Lagos 30        &  30 & -0.00 & Diverse \\
Lagos 50        &  50 & -0.04 & Diverse \\
Lagos 100       & 100 & -0.12 & Diverse \\
Lagos 145       & 145 & +3.06 & Scalarized \\
\midrule
Mexico City 30  &  30 & -0.00 & Diverse \\
Mexico City 50  &  50 & -0.15 & Diverse \\
Mexico City 100 & 100 & -0.26 & Diverse \\
Mexico City 150 & 150 & -0.46 & Diverse \\
\midrule
Paris 30        &  30 & -2.03 & Diverse \\
Paris 50        &  50 & -0.28 & Diverse \\
Paris 100       & 100 & -1.56 & Diverse \\
Paris 150       & 150 & -0.11 & Diverse \\
\midrule
Rio 30          &  30 & -0.03 & Diverse \\
Rio 50          &  50 & -0.03 & Diverse \\
Rio 100         & 100 & -0.12 & Diverse \\
Rio 150         & 150 & +0.02 & Scalarized \\
\midrule
Tokyo Bay 30    &  30 & -0.02 & Diverse \\
Tokyo Bay 50    &  50 & -1.04 & Diverse \\
Tokyo Bay 100   & 100 & -0.09 & Diverse \\
Tokyo Bay 150   & 150 & -0.40 & Diverse \\
\midrule
\multicolumn{2}{l}{\textbf{Average}} & \textbf{-0.18} & \textbf{Diverse} \\
\multicolumn{2}{l}{\textbf{Win--loss}} & \multicolumn{2}{r}{2 wins, 18 losses} \\
\bottomrule
\end{tabular}
\end{table}

The direct comparison is decisive:
\begin{itemize}
  \item Diversity-guided probing wins on 18 of 20 instances.
  \item Scalarization-guided selection wins on only two instances, both
    at the largest size tier.
  \item The average deficit of scalarization-guided selection relative
    to diversity-guided probing is \(0.18\%\).
\end{itemize}

Thus, while scalarization-guided selection is competitive and sometimes
superior on difficult large instances, diversity-guided probing is the
better general-purpose strategy on this benchmark set.

\subsection{Hypervolume-Over-Time Analysis}

The hypervolume-over-time curves confirm the same qualitative picture.
The diversity-guided strategy tends to achieve stronger early coverage
of the approximation set, while the scalarization-guided strategy is
more selective and therefore sometimes slower to spread across the front.
On the hardest instances, however, scalarization-guided selection can
eventually overtake diversity-guided probing by focusing search effort
on strategically important trade-off directions.

\begin{figure}[htbp]
\centering
\includegraphics[width=\textwidth]{../hv_experiment_results_diverse_vs_scalarized_max150/fig_combined.png}
\caption{Normalised hypervolume over time for all benchmark instances in
the Diverse Probe PLS vs Scalarized PLS comparison.}
\label{fig:hv_curves}
\end{figure}

\subsection{Visual Summary}

\Cref{fig:bar_summary} summarises the final HV values across all
instances.

\begin{figure}[htbp]
\centering
\includegraphics[height=0.86\textheight,keepaspectratio]{../hv_experiment_results_diverse_vs_scalarized_max150/fig_barchart.png}
\caption{Final normalised hypervolume by instance for Pure PLS, Diverse
Probe PLS, and Scalarized PLS.}
\label{fig:bar_summary}
\end{figure}

% ══════════════════════════════════════════════════════════════════════
\section{Discussion}\label{sec:discussion}

\paragraph{Why does diversity-guided probing work so well?}
The SIMS objective space is heterogeneous and often exhibits multiple
structurally distinct trade-off regions.  A diversity-guided policy
allocates search effort across these regions more evenly, which is
especially valuable when the archive is already large and the main risk
is oversampling dense regions.  This explains why diversity-guided
probing is the stronger default strategy on most instances.

\paragraph{When does scalarization-guided selection help?}
Scalarization-guided selection is most effective when the search must be
steered toward specific trade-off directions that are difficult to reach
through broad coverage alone.  The Lagos 145 and Rio 150 instances
illustrate this regime: on these large instances, directional archive
steering yields the best final HV.

\paragraph{Small-instance behaviour.}
On the 30-image tier, the differences between the strategies are
minimal.  This suggests that the baseline already explores the search
space sufficiently well under the available budget, leaving little room
for parent-selection improvements.

\paragraph{Medium and large instances.}
From 100 images onward, both enhancements become clearly beneficial over
the baseline.  The gains are particularly strong on Paris 100, Mexico
City 150, Rio 150, and Tokyo Bay 150, indicating that parent-selection
policy matters most when the archive is large and the search budget must
be allocated carefully.

\paragraph{Practical recommendation.}
If a single enhanced strategy must be chosen, the results support
\emph{diversity-guided probing} as the default.  Scalarization-guided
selection remains valuable as a complementary strategy, especially for
large instances where directional exploitation can dominate broad
coverage.

% ══════════════════════════════════════════════════════════════════════
\section{Related Work}\label{sec:related}

Pareto Local Search has long been recognised as a strong local-search
framework for multi-objective combinatorial optimisation.  A central
theme in the literature is that PLS performance depends not only on the
neighbourhood operator, but also on archive management and search
scheduling.

Diversity-oriented archive exploration is closely related to the broader
idea of maintaining spread in multi-objective search, which appears in
evolutionary algorithms, decomposition methods, and archive management
schemes.  The diversity-guided probing strategy studied here adapts this
principle to the parent-selection stage of PLS.

Scalarization-guided search has an equally rich history.  Weighted
Chebycheff and related scalarizations are standard tools for expressing
trade-off directions in multi-objective optimisation.  In the present
context, scalarization is not used to replace Pareto dominance, but to
guide the order in which archive members are selected for local
exploration.

The present study therefore sits at the intersection of two established
ideas---diversity preservation and scalarization-based guidance---and
evaluates them in a controlled PLS setting for the SIMS problem.

% ══════════════════════════════════════════════════════════════════════
\section{Conclusion}\label{sec:conclusion}

We presented a comparative study of two parent-selection enhancements
for Pareto Local Search on the four-objective Satellite Image Mosaic
Selection problem: diversity-guided probing and scalarization-guided
archive selection.

The experiments show that:
\begin{itemize}
  \item both strategies improve over the baseline on medium and large
    instances,
  \item diversity-guided probing is the stronger default strategy,
    improving final HV by \(+3.60\%\) on average over the baseline,
  \item scalarization-guided selection is also effective, improving final
    HV by \(+3.42\%\) on average over the baseline,
  \item direct comparison favours diversity-guided probing on 18 of 20
    instances, although scalarization-guided selection wins on two large
    instances.
\end{itemize}

These results suggest that broad archive coverage is generally more
valuable than directional exploitation in this benchmark set, but that
scalarization-guided steering remains an important complementary tool
for difficult large instances.

\paragraph{Future work.}
Natural next steps include:
\begin{enumerate}
  \item adaptive switching between diversity-guided and
    scalarization-guided parent selection during a run,
  \item hybrid parent-selection policies that combine spread and
    directional preference more tightly,
  \item evaluation on the 200-image instances and under larger time
    budgets,
  \item extension of the same comparison to hybrid exact/heuristic
    configurations.
\end{enumerate}

\paragraph{Acknowledgements.}
Computational experiments were conducted using the open-source SIMS
solver framework with trace-based hypervolume evaluation.

\end{document}
